\documentclass[11pt]{article}

\usepackage{amsmath,amssymb}
\usepackage[margin=1in]{geometry}

\title{Gaussian Limit of the Hyperbolic Gravitational-Wave Likelihood}
\author{}
\date{}

\begin{document}
\maketitle

\section*{Setup}

For one detector and one frequency bin, define the complex residual
\begin{equation}
  r_k = d_k - h_k,
\end{equation}
and write the Bilby per-bin scale as
\begin{equation}
  s_k^2 = \frac{S_n(f_k)\,T}{4},
\end{equation}
where $S_n(f_k)$ is the one-sided noise power spectral density and $T$ is the data duration.

The hyperbolic likelihood used in the implementation treats the real and imaginary parts of
$r_k$ as a two-dimensional real vector. Define the standardized residual
\begin{equation}
  x_k = \frac{r_k}{\sqrt{s_k^2}},
  \qquad
  q_k = |x_k|^2 = \frac{|r_k|^2}{s_k^2}.
\end{equation}

The per-bin symmetric hyperbolic density is
\begin{equation}
  p(x_k \mid \alpha, \delta)
  =
  \frac{\alpha}{2\pi \delta\, K_1(\alpha\delta)}
  K_0\!\left(\alpha \sqrt{\delta^2 + q_k}\right),
\end{equation}
with shape parameters $\alpha > 0$ and $\delta > 0$, and $K_\nu$ the modified Bessel function
of the second kind.

\section*{Asymptotic Limit}

Take the limit
\begin{equation}
  \alpha, \delta \to \infty,
  \qquad
  \lambda = \frac{\alpha}{\delta}
  \quad \text{fixed}.
\end{equation}

Use the large-argument asymptotic expansion
\begin{equation}
  K_\nu(z) \sim \sqrt{\frac{\pi}{2z}} e^{-z}
  \qquad (z \to \infty).
\end{equation}

Applying this to both Bessel functions gives
\begin{align}
  p(x_k \mid \alpha, \delta)
  &\sim
  \frac{\alpha}{2\pi\delta}
  \frac{\sqrt{\frac{\pi}{2\alpha\sqrt{\delta^2 + q_k}}}
        e^{-\alpha\sqrt{\delta^2 + q_k}}}
       {\sqrt{\frac{\pi}{2\alpha\delta}} e^{-\alpha\delta}} \\
  &=
  \frac{\alpha}{2\pi\delta}
  \sqrt{\frac{\delta}{\sqrt{\delta^2 + q_k}}}
  \exp\!\left[-\alpha\left(\sqrt{\delta^2 + q_k} - \delta\right)\right].
\end{align}

Now expand the square root:
\begin{equation}
  \sqrt{\delta^2 + q_k}
  =
  \delta \sqrt{1 + \frac{q_k}{\delta^2}}
  =
  \delta + \frac{q_k}{2\delta} + O(\delta^{-3}).
\end{equation}

Therefore
\begin{equation}
  \alpha \left(\sqrt{\delta^2 + q_k} - \delta\right)
  =
  \frac{\alpha}{2\delta} q_k + O(\delta^{-2})
  \to
  \frac{\lambda q_k}{2},
\end{equation}
and also
\begin{equation}
  \sqrt{\frac{\delta}{\sqrt{\delta^2 + q_k}}} \to 1.
\end{equation}

Hence the limiting density is
\begin{equation}
  p(x_k \mid \alpha, \delta)
  \to
  \frac{\lambda}{2\pi}
  \exp\!\left(-\frac{\lambda q_k}{2}\right).
\end{equation}

This is the density of a two-dimensional zero-mean Gaussian with precision $\lambda$.

\section*{Limit in Terms of the Residual $r_k$}

Since $q_k = |r_k|^2 / s_k^2$, the density in the original residual variable is
\begin{equation}
  p(r_k)
  \to
  \frac{\lambda}{2\pi s_k^2}
  \exp\!\left(-\frac{\lambda |r_k|^2}{2 s_k^2}\right).
\end{equation}

Equivalently, the per-bin log-likelihood becomes
\begin{equation}
  \log p(r_k)
  \to
  \log \lambda - \log(2\pi s_k^2) - \frac{\lambda |r_k|^2}{2 s_k^2}.
\end{equation}

For the special case $\alpha = \delta \to \infty$, one has $\lambda = 1$, so
\begin{equation}
  \log p(r_k)
  \to
  -\log(2\pi s_k^2) - \frac{|r_k|^2}{2 s_k^2}.
\end{equation}

Substituting the Bilby scale
\begin{equation}
  s_k^2 = \frac{S_n(f_k)\,T}{4},
\end{equation}
gives
\begin{equation}
  \log p(r_k)
  \to
  -\log\!\left(\frac{\pi S_n(f_k)\,T}{2}\right)
  - \frac{2|r_k|^2}{T\,S_n(f_k)}.
\end{equation}

This is exactly the standard complex Gaussian per-bin likelihood used by Bilby. Summing over
independent frequency bins and detectors reproduces the usual Gaussian/Whittle gravitational-wave
likelihood, up to the same additive normalization constants already present in the Gaussian model.

\section*{Conclusion}

The symmetric hyperbolic likelihood has a Gaussian limit only when
\begin{equation}
  \alpha, \delta \to \infty
  \quad \text{with} \quad
  \lambda = \frac{\alpha}{\delta}
  \quad \text{fixed}.
\end{equation}
The limiting Gaussian precision is $\lambda$. In particular, choosing $\alpha = \delta$ and
taking both to infinity yields the standard Bilby Gaussian likelihood.

\end{document}
